% Clase 9 --- Continuidad de la componente tangencial (§3.2).
% Circuito rectangular delgado a caballo de la superficie. Como la fuerza es
% conservativa, la circulación es cero; los lados cortos se desprecian y quedan
% sólo los tramos tangenciales: E1t = E2t. Adentro E = 0, luego afuera la
% tangencial también es nula.
\begin{center}
\begin{tikzpicture}[scale=1]

  % conductor
  \fill[figgray!14] (-3.3,-1.6) rectangle (3.3,0);
  \draw[figline, line width=1.2pt] (-3.3,0) -- (3.3,0);
  \node[figlblaux] at (0,-1.1) {conductor: \ $\vec E = 0$};


  % circuito rectangular delgado
  \draw[figred, line width=1.2pt] (-1.3,0.32) rectangle (1.3,-0.32);
  % sentido de circulación
  \draw[figvecres] (-0.5,0.32) -- (0.5,0.32);
  \draw[figvecres] (0.5,-0.32) -- (-0.5,-0.32);
  \node[figlbl, figred, anchor=south] at (0,0.4) {tramo \emph{afuera}};
  \node[figlbl, figred, anchor=north] at (0,-0.4) {tramo \emph{adentro}};

  % los lados cortos, despreciables
  \draw[figaux, line width=0.5pt] (1.3,0) -- (2.3,0.85);
  \node[figlblaux, anchor=west, align=left] at (2.35,0.85)
    {lados cortos:\\ despreciables};

  % componente tangencial de afuera
  \draw[figvec] (-2.9,0.95) -- (-1.85,0.95);
  \node[figlbl, figblue, anchor=south] at (-2.4,1.0) {$E_{1t}$ (afuera)};
  \node[figlblaux, anchor=east] at (-1.45,-0.62) {$E_{2t}=0$};

  % lectura
  \node[figlbl, align=center] at (0,-2.35)
    {$\displaystyle \oint \vec E \cdot d\vec l = 0
      \ \Longrightarrow\ E_{1t} = E_{2t} = 0$};

\end{tikzpicture}\\[0.5em]
{\small\itshape Este es el argumento que \textbf{le faltaba al Resnick}
(Clase 7): la circulación nula ---consecuencia de que la fuerza electrostática
es conservativa--- vale para el rectángulo entero, que está mitad adentro y
mitad afuera. Achatándolo, los lados cortos no aportan y sobreviven sólo las
componentes \textbf{tangenciales}, que resultan iguales. Como adentro el campo
es nulo, afuera la tangencial también lo es: el campo sale perpendicular.}
\end{center}
